\section{Des germes APP au registre de mélange et de masse}
\label{sec:semillas-registro-mezcla-autonoma}
\index{germe!factorisation somme--produit}
\index{mélange!provenance dans les germes}

La loi physique reçoit un domaine qui a déjà conservé l'incompatibilité des
deux évaluations APP. Si
\(\mathcal S_+\) et \(\mathcal S_\times\) désignent les germes orientés des
feuillets additif et multiplicatif, alors
\begin{equation}
 \Omega_{\rm seed}=\mathcal S_+\times\mathcal S_\times,
 \qquad
 |\Omega_{\rm seed}|=324^2=104,976.
 \label{eq:dominio-semillas-md-autonoma}
\end{equation}
La factorisation exacte de l'émetteur produit
\begin{equation}
 104,976
 \longrightarrow18\times26=468\ U_6
 \longrightarrow243\ w_6\subset\mathbb F_3^6,
 \qquad |\mathbb F_3^6|=729.
 \label{eq:embudo-semillas-md-autonoma}
\end{equation}
Les quatre cardinaux comptent des objets distincts : conditions initiales de
deux feuillets, émissions décimales, mots ternaires publiés et espace ambiant
algébrique. Chaque germe conserve, outre son émission, multiplicité,
diagonale, feuillet, orbite, signature d'émission, résidus modulo trois et modulo
neuf, rotations, changements de feuillet, extrémité, déplacements et comptages
cardinaux. Ces données constituent le vocabulaire du registre ultérieur.

La factorisation organise six microfibres de cardinal \(1008\), six plans
de Fano de phase, trente-six fibres de cardinal \(28\) et la bande
\(1944=27\cdot72\), diagonalisée par \((\mathbb Z/3)^3\). En moyennant le
mot nonadique du TPK, les projecteurs conditionnels des deux feuillets
définissent le noyau canonique
\begin{equation}
 K_{\rm cat}=\frac13(I+E_++E_\times),
 \qquad
 \operatorname{Spec}(K_{\rm cat})
 =\left\{1^{[1]},
 \left(\frac23\right)^{[42]},
 \left(\frac13\right)^{[425]}\right\}.
 \label{eq:kernel-canonico-semillas-md-autonoma}
\end{equation}
La phase TPK sélectionne le renouvellement actif et conserve les neuf bandes
spatiales ; la moyenne~\eqref{eq:kernel-canonico-semillas-md-autonoma}
publie l'ombre stochastique du calendrier complet.

\subsection{Bigraduation des germes et \(6\) plans d'incidence}

Dans chacune des six microfibres de cardinalité \(1008\), le graphe
biparti conserve \(54\) sommets additifs et \(56\) multiplicatifs. Les
premiers se décomposent en \(36\) sommets de degré \(24\) et \(18\) de
degré \(8\) ; les seconds ont le degré \(18\). Le décompte des incidences
coïncide pour les deux feuillets :
\begin{equation}
 1008=36\cdot24+18\cdot8=56\cdot18.
 \label{eq:bidegraduacion-microfibra-masas-autonoma}
\end{equation}
Cette égalité fixe la bigraduation de l'émetteur avant sa projection sur
les familles de chemins.

Les six mots
\begin{equation}
 001112,\quad010211,\quad100121,\quad112001,\quad121100,\quad211010
 \label{eq:seis-palabras-microfibra-masas-autonoma}
\end{equation}
forment une orbite de \(D_3\simeq S_3\) et possèdent le profil
\(432+4\cdot144=1008\). Le quotient \(1008=36\cdot28\) organise sept
blocs
\[
 Q=\{B_3,B_6,B_9,P_1,P_2,P_3,f\},
\]
où \((B_3,B_6,B_9)\) sont les tétrades transversales et
\((P_1,P_2,P_3,f)\), les quatre tétrades latérales. Pour l'un des six
mots \(a\), soit
\[
 \mathcal R_a
 :=\{(p,t)\in\mathcal S_+\times\mathcal S_\times:
       U_6(p,t)\text{ publie le mot }a\}
\]
l'ensemble de ses mille huit chemins, et posons
\(R_9:=\mathbb Z/9\mathbb Z\). La carte orientée étant fixée, la projection
de phase a pour type
\begin{equation}
 q_\beta:\mathcal R_a\longrightarrow\binom{R_9}{2},
 \qquad
 q_\beta(p,t)
 =\{x-5\beta(d_t),x+5\beta(d_t)\}\subset\mathbb Z/9\mathbb Z,
 \qquad x=5(i_p-j_p),
 \label{eq:proyeccion-fano-semillas-masas-autonoma}
\end{equation}
avec \((\beta(N),\beta(E),\beta(S),\beta(O))=(1,2,3,4)\). Son codomaine est
l'ensemble des trente-six couples non ordonnés de \(R_9\). Chaque
\(X\in Q\) est une tétrade d'orbites ;
son support de chemins et sa coupe sur une fibre sont définis par
\[
 \operatorname{Supp}_{\rm rutas}(X)
 :=\bigcup_{\mathcal O\in X}\mathcal O,
 \qquad
 X_b:=\operatorname{Supp}_{\rm rutas}(X)\cap q_\beta^{-1}(b),
 \qquad b\in\binom{R_9}{2}.
\]
Dans le canal multiplicatif, l'observable source est
\begin{equation}
 \rho:\mathcal S_\times\longrightarrow R_9,
 \qquad
 \rho(i,j,d)=(i+1)(j+1)\pmod9.
 \label{eq:residuo-producto-fano-semillas-rev3}
\end{equation}
Si \(\pi_\times:\mathcal R_a\to\mathcal S_\times\),
\(\pi_\times(p,t)=t\), est la seconde projection, l'observable qui agit
sur les chemins est
\begin{equation}
 \rho_\times:=\rho\circ\pi_\times:\mathcal R_a\longrightarrow R_9,
 \qquad
 \rho_\times(p,t)=(i_t+1)(j_t+1)\pmod9.
 \label{eq:residuo-producto-rutas-fano-semillas-rev3}
\end{equation}
Dans ce qui suit, \(\rho(X)\) abrège le multiensemble des valeurs de
\(\rho_\times\) sur \(X_b\), qui est indépendant de la fibre \(b\). Ses
profils sur les tétrades sont
\[
 \rho(B_3)=3^4,
 \qquad
 \rho(B_6)=6^4,
 \qquad
 \rho(B_9)=0^4,
\]
et, pour les trois tétrades latérales,
\begin{equation}
 \begin{array}{c|c|c}
 r&\text{signature active de }P_r&R(P_r)\\
 \hline
 1&933&\{5,7\}\\
 2&393&\{4,8\}\\
 3&339&\{1,2\}.
 \end{array}
 \label{eq:perfiles-laterales-fano-semillas-rev3}
\end{equation}
Chaque résidu de \(R(P_r)\) apparaît deux fois parmi les quatre chemins de
\((P_r)_b\).

\begin{theorem}[Morphisme d'incidence des germes et six plans de Fano]
\label{thm:semillas-seis-fano-fase-rev3}
Pour \(\delta\in R_9\), \(h\in\{3,6,0\}\) ---correspondant,
respectivement, à \(B_3,B_6,B_9\)--- et \(r\in\{1,2,3\}\), soit
\begin{equation}
 N_\delta(B_h,P_r)
 :=\sum_{b\in\binom{R_9}{2}}
 \#\left\{(x,y)\in(B_h)_b\times(P_r)_b:
 \rho_\times(y)-\rho_\times(x)=\delta\right\}.
 \label{eq:conteo-incidencia-fano-semillas-rev3}
\end{equation}
Alors
\begin{equation}
 N_\delta(B_h,P_r)=
 \begin{cases}
 288,&h+\delta\in R(P_r),\\
 0,&h+\delta\notin R(P_r).
 \end{cases}
 \label{eq:conteo-288-cero-fano-semillas-rev3}
\end{equation}
Si \(\delta\in R_9^\times\), pour chaque \(B_h\) il existe un unique \(P_r\)
avec comptage \(288\) ; on obtient donc une bijection
\[
 \theta_\delta:\{B_3,B_6,B_9\}
 \xrightarrow{\sim}\{P_1,P_2,P_3\}.
\]
En écrivant l'image de \((B_3,B_6,B_9)\) au moyen de ses indices, les six
bijections sont
\begin{equation}
 \begin{array}{c|c@{\qquad}c|c}
 \delta&\theta_\delta&\delta&\theta_\delta\\
 \hline
 1&(2,1,3)&5&(2,3,1)\\
 2&(1,2,3)&7&(3,2,1)\\
 4&(1,3,2)&8&(3,1,2).
 \end{array}
 \label{eq:tabla-theta-delta-fano-semillas-rev3}
\end{equation}
Soit \(P=\{P_1,P_2,P_3\}\) et, pour \(p\in P\), soit
\[
 \mu(p):=\bigl\{\{f,p\},P\setminus\{p\}\bigr\}
\]
le couplage parfait associé sur les quatre tétrades latérales.
Alors
\begin{equation}
 \mathcal D_\delta
 :=\bigl\{\{B_3,B_6,B_9\}\bigr\}
 \cup
 \bigl\{\{B,q,q'\}:B\in\{B_3,B_6,B_9\},
       \{q,q'\}\in\mu(\theta_\delta(B))\bigr\}
 \label{eq:sistema-fano-semillas-rev3}
\end{equation}
est un système de Steiner \(\operatorname{S}(2,3,7)\). Les six unités
produisent les six plans de Fano sur ces sept blocs qui contiennent la
droite transversale \(\{B_3,B_6,B_9\}\).
\end{theorem}

\begin{proof}
Fixons une fibre \(b\). Dans \((B_h)_b\), il existe quatre choix de
résidu \(h\). Si \(h+\delta\in R(P_r)\), le profil
\eqref{eq:perfiles-laterales-fano-semillas-rev3} fournit exactement deux
choix dans \((P_r)_b\) avec ce résidu ; s'il n'appartient pas au support, il
n'en fournit aucun. Comme il y a trente-six fibres,
\(N_\delta=36\cdot4\cdot2=288\) dans le premier cas et zéro dans le second.
Les trois supports latéraux partitionnent \(R_9^\times\), de sorte qu'une unité
\(\delta\) détermine un unique latéral pour chaque bloc transversal. La lecture
directe de ces supports donne le tableau
\eqref{eq:tabla-theta-delta-fano-semillas-rev3} et épuise les \(3!=6\)
bijections.

La droite transversale couvre les trois couples entre \(B_3,B_6,B_9\). Pour un
\(B\) fixé, les deux arêtes du couplage \(\mu(\theta_\delta(B))\)
couvrent une fois les quatre points latéraux. Les trois couplages sont
distincts et forment la factorisation de \(K_4\) en ses trois couplages
parfaits ; par conséquent, chaque
couple latéral apparaît exactement une fois. Les sept triplets contiennent ainsi
chacun des vingt-et-un couples de points une fois, ce qui est la propriété
\(\operatorname{S}(2,3,7)\). Réciproquement, tout plan de Fano contenant
la droite transversale associe à chaque \(B\) un couplage parfait distinct
des deux autres et retrouve l'une des six bijections précédentes.
\end{proof}

La correspondance d'ensembles \(\delta\mapsto\theta_\delta\) n'est ni une
action ni un homomorphisme : \(R_9^\times\simeq C_6\) est abélien pour le
produit, tandis que \(S_3\) ne l'est pas, et les unités ne forment pas un sous-groupe
additif. Le déplacement \(\delta\) indexe six translations des profils
du produit.

Le même calcul et le même tableau se vérifient dans les six microfibres. Il faut
néanmoins conserver la différence entre support et rôle : le support de
\(f\) coïncide avec l'un des supports latéraux et, par conséquent, l'observable
\(\rho\) ne retrouve pas à lui seul la décomposition \(P\sqcup\{f\}\). Cette
distinction exige la signature additive déjà déclarée dans l'état de germe.

\begin{statusbox}[title={Statut et portée de l'incidence de Fano}]
Le \cref{thm:semillas-seis-fano-fase-rev3} est un théorème fini une fois
fixés \(q_\beta\), l'orientation et l'ordre des chiffres : il définit six jauges
exactes d'incidence sur les blocs. Il organise phase et incidence dans le domaine
des germes ; il ne sélectionne ni particule, ni masse, ni chemin. Sa réalisation physique
s'effectue ensuite au moyen des morphismes du registre enrichi.
\end{statusbox}

Le passage à la physique conserve la chaîne
\begin{equation}
 \begin{aligned}
 \text{germe raffiné}
 &\longrightarrow(\text{signatures }+,\times;U_6;w_6;
                   \text{mémoire de feuillet})\\
 &\longrightarrow\text{registre CT108}
 \longrightarrow(\tau_{12},D_k,\Omega,P_6,\nu_{120},\nu_{270})\\
 &\longrightarrow\text{fibre et multisection}
 \longrightarrow(B_D,B_\Omega)
 \longrightarrow\text{transport CKM/PMNS}\\
 &\longrightarrow\text{caractère et tours}
 \longrightarrow\text{section de masse}.
 \end{aligned}
 \label{eq:cadena-semilla-mezcla-masa-autonoma}
\end{equation}
CKM réalise un changement de base des opérateurs de quark,
\begin{equation}
 B_d^{(u)}=V_{\rm CKM}B_dV_{\rm CKM}^*,
 \qquad
 R_{\rm act}^{B_d^{(u)}}
 =V_{\rm CKM}R_{\rm act}^{B_d}V_{\rm CKM}^*,
 \label{eq:ckm-transporte-semillas-autonoma}
\end{equation}
et préserve spectre et déterminant. Le mélange spécifie la représentation
de l'opérateur dans une base de saveur ; le registre conserve les valeurs propres et leur
généalogie.

\subsection{Lecteur neutrino contenu dans le registre}

Le secteur neutre dispose d'un extracteur entier construit exclusivement
avec des données du registre. Pour chaque bloc dodécaphasique \(E_m\), nous définissons
\begin{equation}
 T(m)=\sum_{t\in E_m}(-1)^{\kappa(t)}\varepsilon(t)
       \mathbf1_{\{d(t)=9\}},
 \qquad
 s_C^{(m)}(\nu)=T(m)-T(m+6),
 \label{eq:lector-neutrino-ledger-autonoma}
\end{equation}
et
\begin{equation}
 \tau_m^\nu=\operatorname{sgn}s_C^{(m)}(\nu).
 \label{eq:trit-neutrino-ledger-autonoma}
\end{equation}
Les quatre triangles de saut quatre produisent
\begin{equation}
 \nu_{120}=\sum_{j=1}^{4}
 \operatorname{sgn}\!\left(\sum_{m\in\mathcal T_j}\tau_m^\nu\right),
 \qquad
 \nu_{270}=\sum_{m=1}^{12}\tau_m^\nu\tau_{m+3}^\nu.
 \label{eq:invariantes-neutrino-ledger-autonoma}
\end{equation}
L'inversion bidirectionnelle transforme
\(\nu_{120}\mapsto-\nu_{120}\) et conserve \(\nu_{270}\). Cette parité
définit l'architecture historique d'un lecteur indépendant des grandeurs
métrologiques utilisées comme cible. Son application
littérale au registre vivant \(81\times108\), en prenant comme
\(\varepsilon(t)\) l'unique bit global de signe conservé par cette
projection, produit
\[
 s_C^{(m)}(\nu)=0,
 \qquad \tau_m^\nu=\nu_{120}=\nu_{270}=0
\]
aux quatre-vingt-une positions. Les événements antipodaux ont été
identifiés par la projection du registre. Le résultat local est conservé
comme lecteur récupéré et, en même temps, comme test de perte d'information :
l'extracteur neutre de
\cref{sec:extractor-neutro-z0-torsion-rev2} agit avant cette
identification et lit le bit bidirectionnel local, le feuillet, la torsion et la
couronne antipodale. Le zéro obtenu après projection ne dégrade pas cette
construction ; il certifie exactement quelle information le lecteur global efface.
Les référentiels
\(F_\ell,F_\nu\), le lecteur~\eqref{eq:invariantes-neutrino-ledger-autonoma}
et la forme agrégée
\begin{equation}
 G_{ab}=\frac1{\sqrt{|L_a||N_b|}}
 \sum_{c\in L_a}\sum_{d\in N_b}K_{\rm registro}(c,d),
 \qquad U_{\rm int}=\operatorname{polar}(G),
 \label{eq:kernel-ledger-pmns-autonoma}
\end{equation}
fixent ainsi les deux extrémités déjà construites du pont PMNS.

L'égalité cardinale entre les cinquante-six sommets multiplicatifs des
germes et les cinquante-six familles CT108 a été soumise à un premier
contrôle indépendant des grandeurs métrologiques utilisées comme cible.
L'application de Hadamard naturelle sur un seul germe n'atteint
que vingt-neuf familles ; en parcourant les symétries diédriques, les signes et les
translations naturelles, la couverture maximale est de trente-cinq. Par conséquent,
l'entrelaceur agit sur le couple de germes enrichis et conserve feuillet,
retenue et mémoire. La matrice d'incidence comparant les émissions partagées
dans ce domaine double devra être de rang cinquante-six, être naturelle sous
orientation et ne pas se factoriser par des compteurs globaux. Ce critère fixe le
domaine exact de la construction suivante.
